\section*{Chapter \ref{ch:g2:caring-for-nature} --- Caring for Nature}

\begin{solution}{exo:g2:caring-for-nature:1}
The surroundings a \omterm{def:g1:living-or-not:living}{living thing} depends on. For example: the air
we breathe, the water we drink, the food we eat --- all from
\omterm{def:g2:caring-for-nature:environment}{nature}.
\end{solution}

\begin{solution}{exo:g2:caring-for-nature:2}
For example: the beetles under the bark, the woodlice beneath the
log, the moss on top, the hedgehog wintering in the hollow.
\end{solution}

\begin{solution}{exo:g2:caring-for-nature:3}
For example: carry all litter home; look instead of collecting,
and never pick all of a clump; put back every turned stone or
log; stay on paths near fragile spots; watch \omterm{def:g1:animals-around-us:animal}{animals} quietly from
a distance and leave their young alone.
\end{solution}

\begin{solution}{exo:g2:caring-for-nature:4}
Because \omterm{def:g2:caring-for-nature:environment}{nature} has no dustbin lorry: what we leave stays --- a
plastic cap for many, many years, long enough to trap or poison
\omterm{def:g1:animals-around-us:animal}{animals} and still be lying there when today's children are
grown-ups.
\end{solution}

\begin{solution}{exo:g2:caring-for-nature:5}
Picked to the last, the clump can make no \omterm{def:g2:from-seed-to-plant:seed}{seeds} --- so no daisies
there next year --- and the bees lose the \omterm{def:g1:plants-around-us:parts}{flowers} they feed from.
\end{solution}

\begin{solution}{exo:g2:caring-for-nature:6}
They rot down slowly into dark, crumbly soil. Worms, woodlice and
\omterm{def:g1:living-or-not:living}{living things} too small to see do the work; at the end we get
compost --- good soil that feeds next year's \omterm{def:g1:plants-around-us:plant}{plants}.
\end{solution}

\begin{solution}{exo:g2:caring-for-nature:7}
For example: a saucer of water in a heat wave, a wild unmowed
corner in the garden, \omterm{def:g2:from-seed-to-plant:seed}{seeds} for birds in the cold, digging a
small pond.
\end{solution}

\begin{solution}{exo:g2:caring-for-nature:8}
Open answer. Plastic wrappers and bottles usually win. Each kind
should have gone home with its owner --- and then into the right
sorting bin.
\end{solution}

\begin{solution}{exo:g2:caring-for-nature:9}
Leave the \omterm{ex:g2:animals-grow-up:frog}{tadpoles} in the pond: it holds everything they need ---
pond water, food, room to grow \omterm{def:g1:my-body:parts}{legs} --- and a jar does not. If
Emma wants to help, she can watch them there through the spring,
keep the bank untrampled, and perhaps help make a new pond
\omterm{def:g2:where-animals-live:habitat}{habitat} like the school one.
\end{solution}

\begin{solution}{exo:g2:caring-for-nature:10}
A million ``nothings'' in the bad direction is a meadow buried in
caps --- small harms multiply. In the good direction, a million
small cares --- litter carried home, taps turned off, waste
sorted --- add up just as surely: \omterm{def:g2:caring-for-nature:environment}{nature} is mostly cared for in
choices exactly that small.
\end{solution}
